% type-theory-analysis.tex — Companion analysis to C4-FLOW-2026
% Flow Links Through the Lens of Type Theory and Substructural Logic
% Working document — not for publication

\documentclass[11pt]{article}
\usepackage[margin=1.2in]{geometry}
\usepackage{amsmath,amssymb,amsthm}
\usepackage{mathpartir}
\usepackage{stmaryrd}
\usepackage{enumitem}
\usepackage{array}
\usepackage{booktabs}
\usepackage{hyperref}

\newtheorem{definition}{Definition}
\newtheorem{proposition}{Proposition}
\newtheorem{observation}{Observation}
\newtheorem{theorem}{Theorem}
\newcommand{\sess}[1]{\mathsf{#1}}
\newcommand{\tp}[1]{\mathsf{#1}}
\newcommand{\kw}[1]{\mathbf{#1}}
\newcommand{\Perms}{\mathsf{Perms}}
\newcommand{\FlowDir}{\mathsf{FlowDir}}
\newcommand{\Sem}{\mathsf{Sem}}
\newcommand{\eventually}{\Diamond}
\newcommand{\always}{\Box}

\title{Flow Links Through the Lens of Type Theory\\and Substructural Logic}
\author{Analysis companion to C4-FLOW-2026}
\date{March 2026}

\begin{document}
\maketitle

\tableofcontents
\bigskip

%% ═══════════════════════════════════════════════════════════════════
\section{Introduction and Summary of Findings}
%% ═══════════════════════════════════════════════════════════════════

This document examines whether the three filesystem link types---hard links,
symbolic links, and flow links---admit a natural characterization in type
theory, and whether such a characterization provides genuine insight beyond
notational reformulation.

\textbf{Summary of findings.} The type-theoretic analysis is not uniformly
productive. Some angles yield real structure; others are notation shopping.
Specifically:

\begin{enumerate}[label=(\roman*)]
  \item \textbf{Substructural types: genuinely illuminating.} The three link
    types correspond precisely to the structural rules of contraction,
    weakening, and a delayed-contraction principle that does not appear in
    standard substructural logics. This is the strongest finding.
  \item \textbf{Session types: genuinely useful.} Flow links have a natural
    session-typed protocol that distinguishes the three directional forms and
    captures the ongoing, stateful interaction between endpoints. This could
    directly inform API design.
  \item \textbf{Modal/temporal logic: moderately illuminating.} The ``pending''
    state of flow links has a clean reading in temporal logic ($\eventually$
    convergence). This is real structure but somewhat obvious once stated.
  \item \textbf{Effect systems: useful for API design.} An effect system can
    track flow relationships and prevent invalid compositions. This is
    engineering value, not deep theory.
  \item \textbf{Dependent types: notation without insight.} The permission
    composition table can be encoded as a dependent function, but the encoding
    does not reveal structure that the table itself does not already show.
  \item \textbf{Ownership/borrowing: partially illuminating.} The Rust analogy
    is imperfect but clarifies one genuine point about aliasing semantics.
\end{enumerate}

We develop each in turn, noting where the formalization is load-bearing and
where it is merely cosmetic.

%% ═══════════════════════════════════════════════════════════════════
\section{Substructural Types}
%% ═══════════════════════════════════════════════════════════════════

\subsection{Background: Structural Rules}

In sequent calculus, three structural rules govern how hypotheses may be used:

\begin{description}[style=nextline]
  \item[Weakening] A hypothesis may be discarded (introduced without being
    used). If $\Gamma \vdash C$, then $\Gamma, A \vdash C$.
  \item[Contraction] A hypothesis may be duplicated (used more than once). If
    $\Gamma, A, A \vdash C$, then $\Gamma, A \vdash C$.
  \item[Exchange] Hypotheses may be reordered. If $\Gamma, A, B, \Delta \vdash
    C$, then $\Gamma, B, A, \Delta \vdash C$.
\end{description}

Substructural logics restrict these rules. \emph{Linear} logic drops both
weakening and contraction: every hypothesis must be used exactly once.
\emph{Affine} logic drops contraction but keeps weakening: every hypothesis may
be used at most once. \emph{Relevant} logic drops weakening but keeps
contraction: every hypothesis must be used at least once.

\subsection{Hard Links as Contraction}

A hard link creates a second name for the same inode. The resource (the
inode's data) can be accessed through either name, and removing one name does
not affect the other. In type-theoretic terms, a hard link is \emph{contraction
applied to a filesystem resource}: the single resource $r$ is duplicated into
two bindings $r_1, r_2$ that share the same underlying value.

\begin{definition}[Hard-link contraction]
If $\Gamma, p : \tp{Inode}(i) \vdash C$, then
$\Gamma, p_1 : \tp{Inode}(i), p_2 : \tp{Inode}(i) \vdash C$ where $p_1$ and
$p_2$ are hard links to inode $i$.
\end{definition}

The key property is that contraction is \emph{free}: there is no cost, no
delay, and no approximation. Both names resolve to exactly the same data at all
times. This is because hard links operate within a single volume---there is no
partition, no network, no propagation delay. Contraction is immediate because
the duplication is purely in the namespace, not in the data.

This also explains the volume-local scope constraint: contraction of a shared
resource is only safe when the two names inhabit the same context (the same
inode table). Cross-volume contraction would require the two inode tables to be
a single coherent resource, which they are not.

\subsection{Symbolic Links as Weakening}

A symbolic link introduces a name that \emph{may not resolve}. The link
exists, but its target may be missing, unmounted, or unreachable. In
substructural terms, the link itself admits \emph{weakening}: the referenced
resource may be absent from the context.

\begin{definition}[Symlink weakening]
A symbolic link $s$ with target path $t$ corresponds to an assumption
$s : \tp{Ref}(t)$ where $t$ may or may not appear in the context $\Gamma$.
If $t \notin \Gamma$, the link is ``dangling''---the reference exists but the
resource does not. This is weakening: the reference was introduced without a
corresponding resource to use.
\end{definition}

More precisely, a symbolic link is a name that \emph{requires} the
infrastructure to resolve it, and that resolution may fail. In the standard
structural rules, weakening allows $\Gamma \vdash C$ to imply $\Gamma, A
\vdash C$---an unused hypothesis. The symlink is the dual: a \emph{used}
reference to a possibly-absent resource. This corresponds to an \emph{option
type}: a symlink has type $\tp{Ref}(t)$, and dereferencing it yields
$\tp{Option}(\tp{Inode})$---either the resolved inode or an error.

The symlink's consistency model (synchronous: resolve or fail) is the
enforcement mechanism for this optional semantics. The link does not
\emph{approximate} the target; it either provides the current state or refuses
to provide anything. This is the CP choice in CAP terms, and in type-theoretic
terms, it is the difference between an option type (definite yes-or-no) and a
delayed type (yes, eventually).

\subsection{Flow Links as Eventual Contraction}

Flow links create a relationship between two paths such that the content at one
path will \emph{eventually} appear at the other. This is neither standard
contraction (which is immediate) nor weakening (which admits permanent
absence). It is something else: a \emph{promise of future contraction}.

\begin{definition}[Eventual contraction]
An outbound flow link from $p_1$ to $p_2$ asserts: if $\Gamma, p_1 :
\tp{Content}(v)$ holds at time $t$, then there exists a time $t' > t$ such
that $\Gamma', p_2 : \tp{Content}(v)$ holds at $t'$, provided no further
updates occur to $p_1$ after $t$.
\end{definition}

This is not a standard structural rule. It is a \emph{temporal} structural
rule: contraction that holds not at the current time but at some future time.
The resource is not immediately available at both names; it is available at the
source name now and at the destination name eventually.

\begin{observation}
The three link types correspond to three regimes of the contraction rule:
\begin{center}
\begin{tabular}{lll}
\toprule
\textbf{Link type} & \textbf{Structural rule} & \textbf{Character} \\
\midrule
Hard link & Contraction & Immediate, exact duplication \\
Symbolic link & Weakening & Optional, may be absent \\
Flow link & Eventual contraction & Delayed, guaranteed arrival \\
\bottomrule
\end{tabular}
\end{center}
\end{observation}

\subsection{Does This Classification Carry Weight?}

This is the critical question. A classification is useful only if it predicts
something or rules something out.

\textbf{What it predicts.} The substructural framing predicts that composition
of link types should respect the structural rules. For example:
\begin{itemize}
  \item A hard link to a flow-linked path should yield ``eventual contraction
    composed with immediate contraction''---the content is immediately available
    at the hard-linked name and eventually propagated through the flow
    link. This is correct.
  \item A symbolic link to a flow-linked path should yield ``optional access to
    eventually-contracted content''---the symlink may dangle, and if it
    resolves, the content behind it may be pending. This is correct and
    describes a real scenario (symlink to an inbound flow directory on an
    unmounted volume).
  \item A flow link to a hard-linked path treats the hard link transparently:
    since all hard links to an inode are equivalent, the flow link propagates
    the inode's content regardless of which name it was written through. This
    falls directly out of contraction being upstream of eventual contraction.
\end{itemize}

\textbf{What it rules out.} The substructural framing suggests that
``immediate contraction across a partition boundary'' is ill-typed---you cannot
have a hard link across volumes, and the type system says this is because
contraction requires the resource to be in the same context. This is not a new
observation (it is the well-known volume-local constraint), but the type system
\emph{derives} it from a general principle rather than stating it as an ad hoc
rule.

\textbf{What it does not do.} The substructural framing does not distinguish
outbound from inbound from bidirectional flow. All three are ``eventual
contraction'' with different directionality. To capture direction, we need
session types (Section~\ref{sec:session}).

\subsection{A Substructural Type System for Links}

We can make the above precise with a small type system. Let $\tp{Path}$ be the
type of filesystem paths, and let $\tp{Res}$ be the type of filesystem
resources (inodes, directory entries, content). The three link types correspond
to three type constructors:

\begin{align}
  \tp{HardLink}(p_1, p_2) &: p_1 =_{\tp{Res}} p_2
    \tag{contraction: same resource, two names} \\
  \tp{SymLink}(s, t) &: s \to \tp{Option}(\tp{Res}(t))
    \tag{weakening: may not resolve} \\
  \tp{FlowLink}(p, q, d) &: p \to_d \eventually \tp{Res}(q)
    \tag{eventual contraction}
\end{align}

where $d \in \{\rightarrow, \leftarrow, \leftrightarrow\}$ is the direction and
$\eventually$ is the temporal ``eventually'' modality (Section~\ref{sec:modal}).
The hard link asserts propositional equality of resources. The symbolic link is
a partial function (may fail). The flow link is a directional eventual mapping.

This type system is minimal but makes the following well-typedness judgments:

\begin{mathpar}
  \inferrule
    {\Gamma \vdash p_1 : \tp{Res}(i) \\ \Gamma \vdash p_2 : \tp{Res}(i) \\
     \text{same volume}(p_1, p_2)}
    {\Gamma \vdash \tp{HardLink}(p_1, p_2)}

  \inferrule
    {\Gamma \vdash s : \tp{Path} \\ t : \tp{Path}}
    {\Gamma \vdash \tp{SymLink}(s, t) : \tp{Option}(\tp{Res}(t))}

  \inferrule
    {\Gamma \vdash p : \tp{Path} \\ q : \tp{RemotePath} \\ d : \FlowDir}
    {\Gamma \vdash \tp{FlowLink}(p, q, d) : p \to_d \eventually \tp{Res}(q)}
\end{mathpar}

The ``same volume'' side condition on hard links is the substructural
constraint that contraction requires a shared context.


%% ═══════════════════════════════════════════════════════════════════
\section{Session Types}\label{sec:session}
%% ═══════════════════════════════════════════════════════════════════

\subsection{Why Session Types}

Session types, introduced by Honda (1993) and developed by Honda, Vasconcelos,
and Kubo (1998), type the protocols of communicating processes. A session type
describes the sequence of messages exchanged between two endpoints, including
their types, directions, and branching structure.

Flow links are a natural candidate for session typing because:
\begin{enumerate}
  \item They have two endpoints (local and remote).
  \item Communication is directional.
  \item The protocol is ongoing (not a single message but a stream of updates).
  \item The endpoints maintain state (current content, pending changes).
  \item The protocol has phases (quiescent, propagating, conflicted).
\end{enumerate}

\subsection{Session Types for the Three Flow Directions}

We define session types for each flow direction. Let $!T$ denote sending a
value of type $T$, $?T$ denote receiving a value of type $T$, $\oplus$ denote
internal choice (the sender chooses), $\&$ denote external choice (the receiver
offers), $\mu X$ denote recursion, and $\kw{end}$ denote session termination.

\begin{definition}[Outbound flow session --- local endpoint]
\begin{align*}
  S_{\text{out}}^{\text{local}} &= \mu X.\; \oplus \{
    \sess{update}: !\tp{Content}.\; ?(\sess{ack} \oplus \sess{nack}).\; X, \\
    &\qquad\quad\;\;\,
    \sess{quiescent}: X, \\
    &\qquad\quad\;\;\,
    \sess{close}: \kw{end}
  \}
\end{align*}
\end{definition}

The local endpoint of an outbound flow link repeatedly chooses to either send an
update (and wait for acknowledgment), remain quiescent, or close the link.

\begin{definition}[Inbound flow session --- local endpoint]
\begin{align*}
  S_{\text{in}}^{\text{local}} &= \mu X.\; \& \{
    \sess{update}: ?\tp{Content}.\; !(\sess{ack} \oplus \sess{nack}).\; X, \\
    &\qquad\quad\;\;\;\,
    \sess{quiescent}: X, \\
    &\qquad\quad\;\;\;\,
    \sess{close}: \kw{end}
  \}
\end{align*}
\end{definition}

The inbound case is the dual: the local endpoint \emph{offers} to receive
updates, and the remote endpoint chooses when to send.

\begin{definition}[Bidirectional flow session --- local endpoint]
\begin{align*}
  S_{\text{bi}}^{\text{local}} &= \mu X.\; \oplus \{
    \sess{send}: !\tp{Content}.\; ?(\sess{ack} \oplus \sess{nack} \oplus \sess{conflict}(\tp{Content})).\; X, \\
    &\qquad\quad\;\;\,
    \sess{recv}: ?\tp{Content}.\; !(\sess{ack} \oplus \sess{nack} \oplus \sess{conflict}(\tp{Content})).\; X, \\
    &\qquad\quad\;\;\,
    \sess{quiescent}: X, \\
    &\qquad\quad\;\;\,
    \sess{close}: \kw{end}
  \}
\end{align*}
\end{definition}

The bidirectional case introduces a $\sess{conflict}$ branch: when both
endpoints have modified content concurrently, the protocol may enter a conflict
resolution phase. This directly reflects the paper's observation that
bidirectional flow admits conflicts while unidirectional flow does not.

\subsection{Duality and the Remote Endpoint}

A key property of session types is \emph{duality}: the remote endpoint's type
is the dual of the local endpoint's type ($!$ becomes $?$, $\oplus$ becomes
$\&$, etc.). This means:

\begin{itemize}
  \item For outbound flow, the remote endpoint has session type
    $\overline{S_{\text{out}}^{\text{local}}}$: it \emph{offers} to receive
    updates, awaits content, and responds with ack/nack.
  \item For inbound flow, the remote endpoint has session type
    $\overline{S_{\text{in}}^{\text{local}}}$: it \emph{chooses} when to send
    updates.
  \item For bidirectional flow, both endpoints share the same session type (it
    is self-dual up to the choice of who initiates each exchange).
\end{itemize}

\begin{observation}
Session type duality captures a key property of flow links: the ``outbound''
endpoint's local session type is the ``inbound'' endpoint's remote session
type, and vice versa. An outbound flow link at system A and an inbound flow
link at system B, pointing at each other, form a well-typed session.
\end{observation}

\subsection{Asynchronous Session Types}

The standard session types above assume synchronous communication. Flow links
are asynchronous---the remote endpoint need not be reachable. Asynchronous
session types (Gay and Vasconcelos, 2010; Mostrous and Vasconcelos, 2011)
address this by allowing sends to be non-blocking and introducing explicit
buffer types.

Under asynchronous session types, the outbound flow session becomes:

\begin{align*}
  S_{\text{out,async}}^{\text{local}} &= \mu X.\; \oplus \{
    \sess{update}: !\tp{Content}.\; X, \\
    &\qquad\quad\;\;\,
    \sess{ack\_recv}: ?(\sess{ack} \oplus \sess{nack}).\; X, \\
    &\qquad\quad\;\;\,
    \sess{quiescent}: X, \\
    &\qquad\quad\;\;\,
    \sess{close}: \kw{end}
  \}
\end{align*}

The send and the acknowledgment are \emph{decoupled}: the local endpoint can
send multiple updates before receiving any acknowledgment. This matches the
reality of flow links, where multiple local modifications may accumulate during
a partition and propagate as a batch when connectivity is restored.

\subsection{What Session Types Buy You}

The session-typed formulation has practical value for API design:

\begin{enumerate}
  \item \textbf{Protocol compliance.} A session-typed API ensures at compile
    time that a flow link client follows the correct protocol. Sending data
    through an inbound-only flow link is a type error.
  \item \textbf{Deadlock freedom.} Well-typed sessions in the standard
    (synchronous) theory are deadlock-free. In the asynchronous case,
    progress guarantees are weaker but still meaningful.
  \item \textbf{Conflict handling.} The bidirectional session type makes
    conflict handling \emph{mandatory}: the $\sess{conflict}$ branch must be
    handled. A client that ignores conflicts is ill-typed. This is a genuine
    safety property.
  \item \textbf{Resource linearity.} Session types are inherently linear (each
    communication action is performed exactly once per iteration). This
    prevents a common class of bugs: using a flow link handle after it has been
    closed, or closing it twice.
\end{enumerate}

\textbf{Verdict.} Session types are the most directly useful formalism for a
language designer implementing flow-aware filesystem APIs. They provide a type
discipline that distinguishes the three directional forms, enforces correct
protocol usage, and makes conflict handling non-optional.


%% ═══════════════════════════════════════════════════════════════════
\section{Effect Systems}\label{sec:effects}
%% ═══════════════════════════════════════════════════════════════════

\subsection{Flow as an Effect}

An effect system (Gifford and Lucassen, 1986) tracks the computational effects
of expressions---memory reads, writes, I/O, exceptions---as part of their
types. Flow links introduce a new kind of effect: \emph{propagation}. Writing
to a flow-linked path is not just a local write; it is a write that will
eventually be visible at a remote location.

We can model this as an effect annotation on filesystem operations:

\begin{align*}
  \tp{write} &: \tp{Path} \to \tp{Content} \to \{\varepsilon\}\; \tp{Unit} \\
  \text{where } \varepsilon &=
    \begin{cases}
      \kw{local} & \text{if path has no flow link} \\
      \kw{flow\_out}(q) & \text{if path has outbound flow to } q \\
      \kw{flow\_in}(q) & \text{if path has inbound flow from } q \\
      \kw{flow\_bi}(q) & \text{if path has bidirectional flow with } q
    \end{cases}
\end{align*}

\subsection{What Effects Buy You}

An effect system for flow links enables the following static checks:

\textbf{Authorization.} Writing to a path with outbound flow implicitly causes
data to leave the system. An effect system can enforce that $\kw{flow\_out}$
effects are only permitted in contexts with appropriate authorization:

\begin{mathpar}
  \inferrule
    {\Gamma \vdash e : \{\kw{flow\_out}(q)\}\; T \\
     \Gamma \vdash \tp{authorized}(q)}
    {\Gamma \vdash e : \{\kw{flow\_out}(q)\}\; T \;\text{ok}}
\end{mathpar}

This captures the security concern from the paper (Section 8.6): outbound flow
links create data exfiltration risk, and the effect system makes this risk
visible in the type.

\textbf{Composition constraints.} Certain effect compositions should be
flagged:
\begin{itemize}
  \item Writing to a path with inbound flow ($\kw{flow\_in}$) may be
    overwritten by the next inbound propagation. The effect system can issue a
    warning: ``local write to inbound-flow path will be superseded.''
  \item Creating an outbound flow link on a path you do not own (no write
    permission) should be rejected: the $\kw{flow\_out}$ effect requires the
    $\kw{write}$ capability on the path.
\end{itemize}

\textbf{Effect polymorphism.} A function that operates on a path without
knowing its flow status can be polymorphic in the flow effect:

\[
  \tp{process\_file} : \forall \varepsilon.\; \tp{Path} \to
    \{\varepsilon\}\; \tp{Result}
\]

This allows generic code to work with both flow-linked and non-flow-linked
paths, while callers that care about flow effects can instantiate $\varepsilon$
and check it.

\textbf{Verdict.} Effect systems provide engineering value---they make flow
effects explicit in types and enable static checking of authorization and
composition. This is useful for API design but does not reveal deep theoretical
structure about flow links themselves. The effect system is a \emph{consumer}
of the flow link semantics, not a \emph{source} of insight about them.


%% ═══════════════════════════════════════════════════════════════════
\section{Modal and Temporal Logic}\label{sec:modal}
%% ═══════════════════════════════════════════════════════════════════

\subsection{The ``Pending'' State as a Modal Assertion}

A flow link that has been declared but not yet fulfilled---because the remote
endpoint is unreachable, or because local changes have not yet propagated---is
in a ``pending'' state. The paper distinguishes this from the ``broken'' state
of a dangling symbolic link: a pending flow link is valid; a dangling symlink
is degenerate.

This distinction maps directly to temporal logic:

\begin{description}
  \item[Symbolic link (dangling):] ``The target exists \emph{now}, or access
    fails.'' No temporal dimension. The assertion $\tp{Res}(t)$ either holds
    in the current state or does not.
  \item[Flow link (pending):] ``The content will \emph{eventually} be
    consistent.'' This is the temporal modality $\eventually$: $\eventually
    (\tp{Res}(p) = \tp{Res}(q))$.
\end{description}

\subsection{Temporal Logic Operators}

In linear temporal logic (LTL), the relevant operators are:

\begin{itemize}
  \item $\eventually \varphi$ (``eventually $\varphi$''): $\varphi$ holds at
    some future time.
  \item $\always \varphi$ (``always $\varphi$''): $\varphi$ holds at all
    future times.
  \item $\varphi \;\mathcal{U}\; \psi$ (``$\varphi$ until $\psi$''):
    $\varphi$ holds until $\psi$ becomes true.
\end{itemize}

The three link types have the following temporal characterizations:

\begin{center}
\begin{tabular}{lp{8cm}}
\toprule
\textbf{Link type} & \textbf{Temporal assertion} \\
\midrule
Hard link &
  $\always (\tp{Res}(p_1) = \tp{Res}(p_2))$ \\
  & Content is identical at all times (invariant). \\
\addlinespace
Symbolic link &
  $\tp{reachable}(t) \Rightarrow \tp{Res}(s) = \tp{Res}(t)$ \\
  & Consistency holds conditionally on reachability. No temporal modality---the
    condition is evaluated at access time. \\
\addlinespace
Flow link (outbound) &
  $\tp{quiescent}(p) \Rightarrow \eventually (\tp{Res}(q) = \tp{Res}(p))$ \\
  & If the source stops changing, the destination eventually converges. \\
\addlinespace
Flow link (bidir.) &
  $\tp{quiescent}(p) \land \tp{quiescent}(q) \Rightarrow \eventually
    (\tp{Res}(p) = \tp{Res}(q))$ \\
  & If both endpoints stop changing, they eventually converge. \\
\bottomrule
\end{tabular}
\end{center}

\subsection{Connection to Modal Type Systems}

Modal type systems (Pfenning and Davies, 2001) internalize the $\always$ and
$\eventually$ modalities as type constructors:

\begin{itemize}
  \item $\always A$ (``necessarily $A$''): a value of type $A$ that is
    available in all contexts (all worlds, all times). In a distributed
    setting, this corresponds to a replicated value.
  \item $\eventually A$ (``possibly $A$''): a value of type $A$ that is
    available in some context (some world, some time). This corresponds to a
    promise or future.
\end{itemize}

Under this reading:
\begin{align}
  \tp{HardLink}(p_1, p_2) &: \always (\tp{Res}(p_1) =_{\tp{Res}} p_2)
    \tag{necessary equality} \\
  \tp{SymLink}(s, t) &: \tp{Res}(s) \to \tp{Option}(\tp{Res}(t))
    \tag{no modality, just partiality} \\
  \tp{FlowLink}(p, q) &: \tp{Res}(p) \to \eventually \tp{Res}(q)
    \tag{eventual delivery}
\end{align}

The flow link's type, $\tp{Res}(p) \to \eventually \tp{Res}(q)$, is
structurally a \emph{promise}: given content at $p$, you will eventually have
content at $q$. This is the standard type of an asynchronous operation in
languages with futures/promises, which suggests that flow links are, at the
type level, \emph{filesystem-level futures}.

\begin{observation}
A flow link from $p$ to $q$ has the same type structure as a future/promise:
it is a value that is not yet available but will eventually become available.
The distinction is that filesystem futures are \emph{persistent} (they survive
reboots, process termination, and system reconfiguration) and
\emph{observable} (they appear in directory listings, metadata queries, and
filesystem dumps). This makes them closer to \emph{persistent promises} in a
distributed system than to in-process futures.
\end{observation}

\textbf{Verdict.} The modal/temporal reading is correct and clarifying. It
provides a precise vocabulary (``eventual contraction,'' ``persistent promise,''
``filesystem-level future'') and connects flow links to well-studied formal
frameworks. But the insight is relatively shallow: the paper already says
``eventual consistency,'' and the modal reading is essentially a
type-theoretic restatement of that phrase. The value is in connecting the
filesystem concept to the broader PL theory landscape, not in revealing
hidden structure.

%% ═══════════════════════════════════════════════════════════════════
\section{Dependent Types and Permission Composition}\label{sec:dependent}
%% ═══════════════════════════════════════════════════════════════════

\subsection{The Permission Composition Table as a Dependent Function}

The paper presents a table mapping (permissions $\times$ flow direction) to
emergent behavior. Can dependent types add structure?

Define:

\begin{align*}
  \FlowDir &= \{\rightarrow, \leftarrow, \leftrightarrow\} \\
  \Perms &= \{r, w, x\}^3 \quad \text{(owner, group, other---simplified to a single rwx triple)} \\
  \Sem &: \FlowDir \times \Perms \to \tp{Behavior}
\end{align*}

where $\tp{Behavior}$ is an enumeration of the emergent semantics:
$\{\text{publishing-point}, \text{staging-area}, \text{collaborative-sync},
  \text{read-cache}, \text{archive-export}, \text{drop-slot},
  \text{blind-relay}, \ldots\}$

As a dependent type:
\[
  \Sem : (d : \FlowDir) \to (p : \Perms) \to \tp{Behavior}(d, p)
\]

This is well-defined but tautological. The dependent type does not \emph{compute}
the behavior---it merely indexes into a lookup table. The type family
$\tp{Behavior}(d, p)$ is no more informative than the table itself. There is
no algebraic structure to exploit: the behaviors do not form a lattice, the
composition is not a homomorphism, and the indexing does not enable proofs
beyond case analysis.

\subsection{Where Dependent Types Could Add Value}

Dependent types become non-trivially useful if we want to \emph{prove
properties} about compositions. For example:

\begin{proposition}[Drop-slot safety]
If $\Perms = \texttt{-wx}$ and $\FlowDir = \rightarrow$, then no read
operation on the local path can succeed, and all written content eventually
appears at the remote path.
\end{proposition}

This is provable in a dependent type theory by case analysis on the permission
bits and the flow direction, combined with the eventual contraction property.
The proof is trivial but machine-checkable, which matters if the permission
composition logic is embedded in a verified filesystem implementation.

\textbf{Verdict.} Dependent types are overkill for understanding the
permission-flow composition but not useless---they become relevant in a
verified implementation where machine-checked proofs of composition properties
are desired. For design-level analysis, the table in the paper is more
informative than the type.


%% ═══════════════════════════════════════════════════════════════════
\section{Ownership, Borrowing, and the Rust Analogy}\label{sec:rust}
%% ═══════════════════════════════════════════════════════════════════

\subsection{Rust's Ownership Model}

Rust distinguishes three relationships between a name and a value:

\begin{description}
  \item[Owned (\texttt{T}):] The name owns the value. When the name goes out
    of scope, the value is dropped. Exactly one owner at a time (affine
    semantics).
  \item[Shared reference (\texttt{\&T}):] Multiple names can read the value
    concurrently. No mutation allowed through shared references.
  \item[Mutable reference (\texttt{\&mut T}):] Exactly one name can read and
    write the value. No other references may coexist.
\end{description}

\subsection{Mapping to Link Types}

\begin{center}
\begin{tabular}{lll}
\toprule
\textbf{Rust concept} & \textbf{Link type} & \textbf{Quality of analogy} \\
\midrule
\texttt{Rc<T>} / shared ownership & Hard link & Good \\
\texttt{\&T} / shared reference & Symbolic link (r-x) & Partial \\
\texttt{\&mut T} / mutable reference & --- & No match \\
Owned \texttt{T} & Single directory entry & Trivial \\
\texttt{Arc<T>} / atomic shared ownership & Flow link (bidir.) & Strained \\
\texttt{Channel<T>} & Flow link (outbound) & Better \\
\bottomrule
\end{tabular}
\end{center}

\textbf{Hard links $\approx$ \texttt{Rc<T>}.} This is the cleanest analogy.
\texttt{Rc<T>} (reference-counted pointer) allows multiple owners of the same
heap allocation, with the allocation freed when the last owner drops it. Hard
links allow multiple directory entries for the same inode, with the inode freed
when the last link is removed. The reference count is literally the inode's
link count. The analogy is almost isomorphic.

\textbf{Symbolic links $\approx$ \texttt{\&T}? No.} A symbolic link is not
a borrow---it does not prevent the target from being removed (there is no
borrow checker enforcing lifetime constraints). A dangling symlink is a
use-after-free in Rust terms, except it is allowed and even common. The analogy
breaks because Rust's references are \emph{checked} (the borrow checker
prevents dangling) while symlinks are \emph{unchecked} (dangling is a runtime
condition, not a compile-time error). A better analogy is \texttt{Weak<T>}: a
weak reference that may or may not still point to a live allocation, and that
must be checked (\texttt{upgrade()}) before use.

\begin{observation}
Symbolic links are \texttt{Weak<T>}, not \texttt{\&T}. Both are references
that may dangle, require runtime checking before use, and do not prevent the
referent from being deallocated.
\end{observation}

\textbf{Flow links $\approx$ channels, not references.} Flow links do not fit
the ownership/borrowing model because they are not about aliasing a single
resource. They are about \emph{transferring} content from one location to
another. This is a channel, not a reference:

\begin{itemize}
  \item An outbound flow link is a \texttt{Sender<Content>}: the local
    endpoint sends content that arrives at the remote endpoint.
  \item An inbound flow link is a \texttt{Receiver<Content>}: the local
    endpoint receives content from the remote endpoint.
  \item A bidirectional flow link is a pair
    \texttt{(Sender<Content>, Receiver<Content>)}.
\end{itemize}

This is a more honest analogy than trying to force flow links into the
reference/borrow framework. Flow links are communication primitives with
persistence, not aliasing primitives.

\textbf{Verdict.} The Rust analogy is useful for exactly one insight: hard
links are \texttt{Rc<T>} and symbolic links are \texttt{Weak<T>}, which
correctly captures the aliasing and dangling properties. For flow links, the
ownership model is the wrong lens entirely; the channel model is more apt.
This confirms the session-type analysis: flow links are fundamentally about
\emph{communication}, not \emph{aliasing}.


%% ═══════════════════════════════════════════════════════════════════
\section{Synthesis: What the Formalisms Reveal}\label{sec:synthesis}
%% ═══════════════════════════════════════════════════════════════════

\subsection{The Two Natures of Filesystem Links}

The type-theoretic analysis reveals that filesystem links have two distinct
natures, and the three link types sit at different points on each axis:

\begin{description}
  \item[Aliasing axis] (substructural types): How many names can refer to the
    same resource, and with what constraints?
    \begin{itemize}
      \item Hard link: unrestricted aliasing (contraction)
      \item Symbolic link: conditional aliasing (weakening---the alias may be
        vacuous)
      \item Flow link: \emph{not aliasing at all}---two distinct resources that
        converge over time
    \end{itemize}
  \item[Communication axis] (session types): What protocol governs the
    interaction between the link's endpoints?
    \begin{itemize}
      \item Hard link: no protocol (same resource, no communication needed)
      \item Symbolic link: request/response (resolve path, get result or error)
      \item Flow link: ongoing asynchronous session (stream of updates,
        acknowledgments, conflict resolution)
    \end{itemize}
\end{description}

\begin{observation}
Hard links are pure aliasing (contraction, no communication). Symbolic links
are conditional aliasing (weakening, synchronous communication). Flow links
are pure communication (no aliasing, asynchronous protocol). The three link
types span from ``all aliasing, no communication'' to ``no aliasing, all
communication.''
\end{observation}

This is the deepest structural finding of the analysis. It means that
``flow link'' is a somewhat misleading name from the aliasing perspective---it
is not a ``link'' in the sense of creating another name for the same thing. It
is a \emph{channel} that happens to be declared in the filesystem namespace.
The paper recognizes this when it says a flow link is ``a declaration of
relationship, not a mechanism of fulfillment,'' but the type theory makes it
precise: a flow link has the type of a communication channel, not the type of
a reference.

\subsection{Implications for Language Design}

For a language designer implementing flow-aware filesystem APIs, the analysis
suggests:

\begin{enumerate}
  \item \textbf{Use session types for the flow link API.} The three directional
    forms should be distinct types with session-typed protocols. An outbound
    flow handle should support only send operations; an inbound flow handle
    should support only receive operations. The bidirectional handle should
    require conflict handling.

  \item \textbf{Track flow effects.} Writes to flow-linked paths should carry
    an effect annotation indicating that the write has propagation
    consequences. This enables static checking of authorization and prevents
    accidental exfiltration.

  \item \textbf{Do not try to unify hard links, symlinks, and flow links under
    a single ``link'' abstraction.} The type-theoretic analysis shows they have
    fundamentally different structure: contraction, weakening, and
    communication. A unified ``Link'' type that covers all three would need to
    be so abstract as to be useless. Instead, keep them as three distinct
    abstractions with shared vocabulary (path, target) but different
    operations and guarantees.

  \item \textbf{Model the ``pending'' state explicitly.} A flow-linked path
    should have an observable state that includes not just the current content
    but the convergence status: synchronized, pending (with metadata about what
    is pending), or conflicted. This is the $\eventually$ modality made
    programmatically observable.

  \item \textbf{Consider linear/affine types for flow link handles.} A flow
    link handle should not be silently discarded (that would leave a flow link
    without a managing process) or duplicated (that would create ambiguity
    about which process is responsible for propagation). Affine or linear
    types enforce this discipline.
\end{enumerate}

\subsection{A Sketch of a Flow-Aware Type System}

Combining the above, a minimal flow-aware type system for filesystem operations
might look like:

\begin{align*}
  \tp{Path} &: \tp{Type} \\
  \tp{Content} &: \tp{Type} \\
  \tp{FlowHandle}[d] &: \tp{Type} \quad \text{(linear, parameterized by direction } d \text{)} \\[6pt]
  \tp{read} &: \tp{Path} \to \{\kw{read}\}\; \tp{Option}(\tp{Content}) \\
  \tp{write} &: \tp{Path} \to \tp{Content} \to \{\kw{write} \cup \varepsilon_{\text{flow}}\}\; \tp{Unit} \\[6pt]
  \tp{flow\_out} &: \tp{Path} \to \tp{RemotePath} \to
    \{\kw{create\_flow}\}\; \tp{FlowHandle}[\rightarrow] \\
  \tp{flow\_in} &: \tp{Path} \to \tp{RemotePath} \to
    \{\kw{create\_flow}\}\; \tp{FlowHandle}[\leftarrow] \\
  \tp{flow\_bi} &: \tp{Path} \to \tp{RemotePath} \to
    \{\kw{create\_flow}\}\; \tp{FlowHandle}[\leftrightarrow] \\[6pt]
  \tp{close\_flow} &: \tp{FlowHandle}[d] \to \{\kw{close\_flow}\}\; \tp{Unit}
    \quad \text{(consumes the linear handle)}
\end{align*}

The effect $\varepsilon_{\text{flow}}$ on $\tp{write}$ is determined by the
flow status of the path: $\kw{local}$ if no flow link exists, $\kw{flow\_out}(q)$
if an outbound flow link to $q$ exists, etc. The $\tp{FlowHandle}$ is linear,
ensuring it must be either closed or transferred, never silently dropped.


%% ═══════════════════════════════════════════════════════════════════
\section{Limitations and Non-Findings}
%% ═══════════════════════════════════════════════════════════════════

Several initially promising directions did not yield substantive results:

\textbf{Homotopy type theory.} One might ask whether the three link types
correspond to different notions of path/equality in HoTT (propositional
equality, paths in a higher groupoid, etc.). We found no meaningful
correspondence. HoTT paths are about identity of types, not about filesystem
naming, and the analogy does not survive scrutiny.

\textbf{Category theory.} Links can be modeled as morphisms in a category of
filesystem states, but this does not add insight beyond what the direct
type-theoretic analysis provides. Hard links are isomorphisms (symmetric,
invertible). Symbolic links are partial morphisms. Flow links are morphisms
in a Kleisli category for the $\eventually$ monad. These are all correct
but do not tell a language designer anything they do not already know.

\textbf{Game semantics.} The interactive nature of flow links suggests a game-
semantic reading (local and remote as player and opponent), but the games are
too simple---the strategies are essentially the session types we already
derived, and game semantics adds overhead without additional precision.

\textbf{Separation logic.} One might model flow links using separation logic's
$*$ (separating conjunction) to reason about disjoint heaps at the two
endpoints. This is viable but reduces to the distributed separation logic
literature, which is well-developed but does not specifically illuminate flow
links as opposed to any other distributed data structure.

In all these cases, the formalism can \emph{express} flow links but does not
\emph{illuminate} them. The two formalisms that provide genuine
illumination---substructural types for the aliasing/structural rule
classification, and session types for the communication protocol---remain the
most productive lenses.


%% ═══════════════════════════════════════════════════════════════════
\section{Conclusion}
%% ═══════════════════════════════════════════════════════════════════

The type-theoretic analysis of flow links yields two principal insights and
several secondary ones.

The first principal insight is \textbf{substructural}: the three filesystem
link types correspond to three structural rules---contraction (hard links),
weakening (symbolic links), and eventual contraction (flow links). The third
is genuinely new; it does not appear in standard substructural logics and
arises from the combination of the contraction rule with a temporal modality.
This classification predicts the composition behavior of link types and
derives the volume-local constraint on hard links from a general principle
rather than an ad hoc rule.

The second principal insight is \textbf{communicative}: flow links are not
links in the aliasing sense but channels in the session-type sense. They sit
at the opposite end of a spectrum from hard links: hard links are pure
aliasing without communication; flow links are pure communication without
aliasing. Session types provide the correct formalism for their directional
forms, naturally capture the distinction between unidirectional and
bidirectional flow, and make conflict handling in the bidirectional case a
type-level obligation.

Together, these two insights characterize flow links as occupying a genuinely
novel position in the type-theoretic landscape: \emph{persistent, filesystem-
level asynchronous channels with eventual-contraction semantics}. This is not
merely a restatement of the paper's informal description; it connects flow
links to specific, well-studied formal frameworks (substructural logic, session
types, temporal modalities) and identifies precisely where existing frameworks
are sufficient and where flow links require extension (the eventual contraction
rule).

For a language designer, the actionable recommendations are: type the three
directional forms as distinct session types; track propagation as an effect;
use linear types for flow link handles; and resist the temptation to unify all
three link types under a single abstraction, because the type theory says they
are fundamentally different kinds of things.

\end{document}
